\chapter{سلاسل ديريشليه والمنتجات الأويلرية}

\section{نطاق الفصل وحالته}

يبني هذا الفصل الجسر التحليلي بين الدوال الحسابية المدروسة في الفصل
السابق والدوال العقدية التي ستقود إلى دالة زيتا ودوال \(L\). ونقطة البدء هي
تحويل متتالية المعاملات \(a(n)\) إلى دالة في المتغير العقدي
\[
D_a(s)=\sum_{n=1}^{\infty}\frac{a(n)}{n^s}.
\]

حالة الفصل في الإصدار \(0.10.0\) هي \textenglish{REVIEWED}: طوبقت
نتائجه وبراهينه مع سجل المعرفات، ونجحت فحوص الجودة والبناء، واكتملت
المراجعة البشرية المستقلة التي أعقبها دمج \textenglish{PR \#4}. وتبقى
توسعة الحلول المختارة ومعالجة الديون التحريرية مطلوبتين قبل بوابة
\textenglish{RELEASE-READY}.

\subsection{المتطلبات السابقة}

يفترض الفصل معرفة:

\begin{itemize}
  \item رموز الرتبة والجمع الجزئي لأبيل من الفصل الثاني.
  \item التقارب المحلي المنتظم والهولومورفية من الفصل الثالث.
  \item الدوال الضربية والالتفاف الديريشلي من الفصل الرابع.
\end{itemize}

\subsection{نواتج التعلم}

بعد إتمام الفصل ينبغي أن يستطيع القارئ:

\begin{enumerate}
  \item التمييز بين التقارب العادي والمطلق والمحلي المنتظم لسلسلة ديريشليه.
  \item تفسير نصفي مستوى التقارب والتقارب المطلق.
  \item تبرير هولومورفية مجموع السلسلة في نصف المستوى المناسب.
  \item تحويل تقدير للمجاميع الجزئية إلى تمثيل تكاملي للسلسلة.
  \item ربط ضرب سلاسل ديريشليه بالالتفاف الديريشلي.
  \item اشتقاق المنتج الأويلري من الضربية والتقارب المطلق.
  \item استرجاع معاملات سلسلة ديريشليه من مجموعها.
  \item اشتقاق المتسلسلات الأساسية المرتبطة بـ\(\mu\)، و\(\Lambda\)،
  و\(\varphi\)، و\(\lambda\).
\end{enumerate}

\section{الترددات الضربية}

نكتب دائمًا
\[
s=\sigma+it,
\qquad \sigma,t\in\mathbb{R}.
\]
ولأن \(n>0\)، نعرف
\[
n^{-s}=e^{-s\log n}=n^{-\sigma}e^{-it\log n}.
\]
إذن
\[
|n^{-s}|=n^{-\sigma}.
\]

يفصل هذا التحليل بين وظيفتين:

\begin{itemize}
  \item الجزء الحقيقي \(\sigma\) يضبط الحجم والتقارب.
  \item الجزء التخيلي \(t\) يولد تذبذبًا بتردد \(\log n\).
\end{itemize}

ولهذا تعد سلاسل ديريشليه نظيرًا ضربيًا لمتسلسلات فورييه: الترددات هنا
ليست الأعداد الصحيحة \(n\)، بل اللوغاريتمات \(\log n\).

\section{التعريف والأمثلة الأولى}

\begin{definition}
سلسلة ديريشليه العامة هي متسلسلة من الشكل
\[
D_a(s)=\sum_{n=1}^{\infty}a(n)n^{-s},
\]
حيث \(a:\mathbb{N}\to\mathbb{C}\) دالة حسابية.
\end{definition}

\begin{example}[دالة زيتا في مجال السلسلة]
إذا كان \(a(n)=1\)، فإن
\[
D_a(s)=\zeta(s)=\sum_{n=1}^{\infty}\frac1{n^s},
\qquad \sigma>1.
\]
هذا التعريف بالسلسلة لا يشمل بعد الاستمرار التحليلي إلى خارج النصف
\(\sigma>1\).
\end{example}

\begin{example}[دالة إيتا]
إذا كان \(a(n)=(-1)^{n-1}\)، فإن
\[
\eta(s)=\sum_{n=1}^{\infty}\frac{(-1)^{n-1}}{n^s}.
\]
تتقارب هذه السلسلة عندما \(\sigma>0\)، بينما تتقارب مطلقًا فقط عندما
\(\sigma>1\). وهذا أول تنبيه إلى أن نصفي مستوى التقارب العادي والمطلق قد
يختلفان.
\end{example}

\begin{example}[سلسلة ذات دعم على المربعات]
إذا كانت \(1_{\square}(n)\) دالة مؤشر المربعات الكاملة، فإن
\[
\sum_{n=1}^{\infty}\frac{1_{\square}(n)}{n^s}
=
\sum_{m=1}^{\infty}\frac1{m^{2s}}
=
\zeta(2s),
\qquad \sigma>\frac12.
\]
\end{example}

\section{أنماط التقارب}

نميز بين ثلاثة أنماط:

\begin{enumerate}
  \item التقارب النقطي عند قيمة محددة لـ\(s\).
  \item التقارب المطلق إذا
  \[
  \sum_{n=1}^{\infty}|a(n)|n^{-\sigma}<\infty.
  \]
  \item التقارب المحلي المنتظم على مجال إذا كان التقارب منتظمًا على كل
  مجموعة مدمجة داخله.
\end{enumerate}

التقارب المحلي المنتظم هو النمط الذي يسمح بنقل الهولومورفية من المجاميع
الجزئية إلى المجموع النهائي. أما التقارب النقطي وحده فلا يكفي لذلك.

\section{الانتقال إلى يمين نقطة التقارب}

\begin{theorem}[نصف مستوى التقارب والهولومورفية]
\label{thm:dirichlet-series-right-half-plane}
\resultid{ANT-THM-05-01}
\provedhere
إذا تقاربت السلسلة
\[
\sum_{n=1}^{\infty}a(n)n^{-s}
\]
عند \(s=s_0\)، فإنها تتقارب محليًا بانتظام في نصف المستوى
\[
\Re(s)>\Re(s_0).
\]
ومن ثم يكون مجموعها هولومورفيًا هناك.
\end{theorem}

\begin{proof}
ضع
\[
b_n=a(n)n^{-s_0},
\qquad
B(x)=\sum_{n\leq x}b_n.
\]
تقارب \(\sum b_n\) يقتضي أن \(B(x)=O(1)\). وإذا كتبنا
\[
w=s-s_0,
\]
فإن الجمع الجزئي يعطي، لأي \(N>M\)،
\[
\sum_{M<n\leq N}b_n n^{-w}
=
B(N)N^{-w}-B(M)M^{-w}
+w\int_M^N B(x)x^{-w-1}\,dx,
\]
مع تعديل طرفي غير مؤثر إذا لم يكن \(M\) عددًا صحيحًا.

لتكن \(K\) مجموعة مدمجة داخل \(\Re(w)>0\). يوجد \(\delta>0\) و\(C_K>0\)
بحيث
\[
\Re(w)\geq\delta,
\qquad |w|\leq C_K
\]
لكل \(w\in K\). وبما أن \(B\) محدودة، فإن ذيل السلسلة محكوم بكمية من
رتبة
\[
O_K(M^{-\delta}).
\]
إذن تحقق السلسلة معيار كوشي بانتظام على \(K\)، فتتقارب محليًا بانتظام.
وكل مجموع جزئي هولومورفي، ولذلك يكون المجموع هولومورفيًا على نصف المستوى.
\end{proof}

\begin{remark}
لا تقول المبرهنة إن التقارب منتظم على نصف المستوى المفتوح كله؛ فالعبارة
الصحيحة اللازمة للتحليل العقدي هي الانتظام على كل مجموعة مدمجة داخله.
\end{remark}

\section{فاصلا التقارب}

\begin{definition}
نعرف فاصل التقارب، في المستقيم الحقيقي الموسع، بـ
\[
\sigma_c
=
\inf\left\{u\in\mathbb{R}:
\text{تتقارب السلسلة لكل }s\text{ ذي }\Re(s)>u
\right\},
\]
مع الاصطلاح \(\inf\varnothing=+\infty\). ويعرف فاصل التقارب المطلق
\(\sigma_a\) باستبدال التقارب بالتقارب المطلق في المجموعة نفسها.
\end{definition}

قد تكون إحدى القيمتين \(-\infty\) أو \(+\infty\)، ولا يقرر التعريف
وحده سلوك السلسلة على خط الحد. وفي الحالات غير المنحلة يكون لدينا نصف
مستوى للتقارب ونصف مستوى التقارب المطلق المحتوى فيه.

\begin{proposition}[الفجوة بين فاصلي التقارب]
\label{prop:dirichlet-series-abscissae-gap}
\resultid{ANT-PROP-05-01}
\provedhere
لكل سلسلة ديريشليه عامة:
\[
\sigma_c\leq\sigma_a\leq\sigma_c+1.
\]
\end{proposition}

\begin{proof}
المتباينة الأولى مباشرة، لأن التقارب المطلق يستلزم التقارب. إذا كان
\(\sigma_c=+\infty\) فلا يبقى ما يثبت. وإذا كان \(\sigma_c=-\infty\)،
فيمكن في الحجة الآتية اختيار \(u\) سالبًا كيفما نشاء، فتنتج
\(\sigma_a=-\infty\). لذا نعرض الحجة الأساسية للحالة
\(\sigma_c\in\mathbb{R}\).

اختر عددًا حقيقيًا \(u>\sigma_c\). تتقارب
\[
\sum_{n=1}^{\infty}a(n)n^{-u},
\]
ومن ثم تكون حدودها محدودة: يوجد \(C>0\) بحيث
\[
|a(n)|n^{-u}\leq C.
\]
إذا كان \(\sigma>u+1\)، فإن
\[
\sum_{n=1}^{\infty}|a(n)|n^{-\sigma}
\leq
C\sum_{n=1}^{\infty}n^{-(\sigma-u)},
\]
والسلسلة في الطرف الأيمن متقاربة. لذلك \(\sigma_a\leq u+1\). وبترك
\(u\downarrow\sigma_c\) نحصل على النتيجة.
\end{proof}

\begin{example}
لدالة إيتا:
\[
\sigma_c=0,
\qquad
\sigma_a=1.
\]
إذن يمكن أن تبلغ الفجوة الحد الأقصى \(1\).
\end{example}

\section{المجاميع الجزئية والتمثيل التكاملي}

لتكن
\[
A(x)=\sum_{n\leq x}a(n).
\]
تربط صيغة أبيل بين نمو \(A(x)\) وتقارب سلسلة ديريشليه.

\begin{theorem}[تمثيل أبيل لسلسلة ديريشليه]
\label{thm:dirichlet-series-abel-integral}
\resultid{ANT-THM-05-02}
\provedhere
إذا كان
\[
A(x)=O(x^\theta)
\]
لبعض \(\theta\in\mathbb{R}\)، فإن السلسلة تتقارب عندما
\(\Re(s)>\theta\)، ولدينا
\[
\sum_{n=1}^{\infty}\frac{a(n)}{n^s}
=
s\int_1^{\infty}A(x)x^{-s-1}\,dx
\]
في هذا النصف من المستوى.
\end{theorem}

\begin{proof}
تطبيق الجمع الجزئي على \(f(x)=x^{-s}\) يعطي
\[
\sum_{n\leq X}a(n)n^{-s}
=
A(X)X^{-s}
+s\int_1^X A(x)x^{-s-1}\,dx.
\]
إذا كانت \(\sigma=\Re(s)>\theta\)، فإن
\[
A(X)X^{-s}=O(X^{\theta-\sigma})\longrightarrow0.
\]
كما أن التكامل يتقارب مطلقًا لأن
\[
|A(x)x^{-s-1}|
\ll x^{\theta-\sigma-1},
\]
وأس القوة أصغر من \(-1\). نترك \(X\to\infty\) فنحصل على الصيغة.
\end{proof}

\begin{example}[الإلغاء يوسع مجال التقارب]
للمعاملات \(a(n)=(-1)^{n-1}\) تكون المجاميع الجزئية محدودة، أي
\(A(x)=O(1)\). لذلك تتقارب \(\eta(s)\) عندما \(\sigma>0\). أما سلسلة
القيم المطلقة فهي \(\zeta(\sigma)\)، فلا تتقارب إلا عندما \(\sigma>1\).
\end{example}

\begin{remark}
المبدأ البنيوي هنا هو أن صغر المجاميع الجزئية، أي وجود إلغاء بين
المعاملات، يدفع نصف مستوى التقارب إلى اليسار.
\end{remark}

\section{الاشتقاق حدًا حدًا}

إذا كانت السلسلة متقاربة مطلقًا في \(\Re(s)>\sigma_a\)، فإنها قابلة
للاشتقاق حدًا حدًا داخل هذا النصف:
\[
D_a'(s)
=
-\sum_{n=1}^{\infty}\frac{a(n)\log n}{n^s}.
\]

والسبب أن كل مجموعة مدمجة تقع على مسافة موجبة من الحد
\(\Re(s)=\sigma_a\). نختار \(\varepsilon>0\) أصغر من تلك المسافة، ثم
نستخدم
\[
\log n\ll_\varepsilon n^\varepsilon
\]
للحصول على تقارب مطلق منتظم لسلسلة المشتقات.

\section{فرادة معاملات سلسلة ديريشليه}

\begin{theorem}[مبرهنة الفرادة]
\label{thm:dirichlet-series-uniqueness}
\resultid{ANT-THM-05-03}
\provedhere
لتكن
\[
F(s)=\sum_{n=1}^{\infty}a(n)n^{-s},
\qquad
G(s)=\sum_{n=1}^{\infty}b(n)n^{-s}
\]
متقاربتين مطلقًا في نصف مستوى مشترك. إذا كان \(F(s)=G(s)\) هناك، فإن
\[
a(n)=b(n)
\]
لكل \(n\geq1\).
\end{theorem}

\begin{proof}
ضع \(c(n)=a(n)-b(n)\). إذا لم تكن جميع المعاملات صفرًا، فليكن \(m\)
أصغر عدد يحقق \(c(m)\neq0\). على المحور الحقيقي، ولـ\(\sigma\) كبيرة، لدينا
\[
0=m^\sigma\sum_{n=m}^{\infty}c(n)n^{-\sigma}
=
c(m)+
\sum_{n>m}c(n)\left(\frac{m}{n}\right)^\sigma.
\]
اختر \(\sigma_0\) في مجال التقارب المطلق. عندما \(\sigma\geq\sigma_0\)،
يمكن السيطرة على الذيل بمتسلسلة متقاربة مشتقة من
\(\sum |c(n)|n^{-\sigma_0}\)، بينما يؤول كل عامل
\((m/n)^{\sigma-\sigma_0}\) إلى الصفر. بالتقارب المسيطر يؤول الذيل إلى
الصفر، فنحصل على \(c(m)=0\)، وهذا تناقض.
\end{proof}

\begin{remark}
تعني الفرادة أن هوية بين سلاسل ديريشليه في نصف مستوى مطلق التقارب ليست
مجرد مساواة دوال؛ إنها هوية معامل بمعامل.
\end{remark}

\section{ضرب السلاسل والالتفاف}

أثبت الفصل الرابع، تحت التقارب المطلق، أن
\[
\left(\sum_{n=1}^{\infty}\frac{f(n)}{n^s}\right)
\left(\sum_{n=1}^{\infty}\frac{g(n)}{n^s}\right)
=
\sum_{n=1}^{\infty}\frac{(f*g)(n)}{n^s}.
\]

هذه الصيغة تجعل تحويل ديريشليه
\[
f\longmapsto D_f(s)
\]
يحول الالتفاف إلى ضرب. لكنها لا تجيز ضرب سلسلتين خارج مجال يبرر إعادة
ترتيب المجموع المزدوج.

\begin{example}
بما أن \(1*1=\tau\)، فإن
\[
\sum_{n=1}^{\infty}\frac{\tau(n)}{n^s}
=
\zeta(s)^2,
\qquad \sigma>1.
\]
\end{example}

\section{المنتج الأويلري العام}

\begin{theorem}[المنتج الأويلري للدالة الضربية]
\label{thm:multiplicative-dirichlet-series-euler-product}
\resultid{ANT-THM-05-04}
\provedhere
لتكن \(f\) دالة ضربية، ولنفترض أن
\[
\sum_{n=1}^{\infty}\frac{|f(n)|}{n^\sigma}<\infty.
\]
عندئذ، لكل \(s\) يحقق \(\Re(s)\geq\sigma\)،
\[
\sum_{n=1}^{\infty}\frac{f(n)}{n^s}
=
\prod_p
\left(
\sum_{k=0}^{\infty}\frac{f(p^k)}{p^{ks}}
\right),
\]
حيث يعرف الجداء بوصفه نهاية الجداءات المنتهية المرتبة بحسب \(p\leq y\).
\end{theorem}

\begin{proof}
لكل \(y\) منتهٍ، يؤدي نشر الجداء إلى
\[
\prod_{p\leq y}
\left(
\sum_{k=0}^{\infty}\frac{f(p^k)}{p^{ks}}
\right)
=
\sum_{\substack{n\geq1\\p\mid n\Rightarrow p\leq y}}
\frac{f(n)}{n^s}.
\]
وحيدية التحليل إلى عوامل أولية تضمن ظهور كل عدد مسموح به مرة واحدة،
وضربية \(f\) تحول حاصل ضرب العوامل المحلية إلى \(f(n)\).

أما تبرير النشر والحد، فيأتي من التقارب المطلق. فإذا كان
\(\Re(s)\geq\sigma\)، فإن \(n^{-\Re(s)}\leq n^{-\sigma}\)، فالفرق بين السلسلة الكاملة
والمجموع المقيد محكوم بذيل من
\[
\sum_{n=1}^{\infty}\frac{|f(n)|}{n^\sigma},
\]
ويؤول هذا الذيل إلى الصفر عندما \(y\to\infty\). لذلك تؤول الجداءات
المنتهية إلى سلسلة ديريشليه المطلوبة.
\end{proof}

إذا كانت \(f\) ضربية تمامًا، فإن
\[
f(p^k)=f(p)^k,
\]
ومن ثم يكون العامل المحلي متسلسلة هندسية:
\[
\sum_{k=0}^{\infty}\frac{f(p^k)}{p^{ks}}
=
\left(1-\frac{f(p)}{p^s}\right)^{-1}
\]
داخل مجال التقارب المطلق.

\section{حاصل ضرب أويلر لدالة زيتا}

بتطبيق المبرهنة على الدالة الثابتة \(1\) نحصل على
\[
\zeta(s)
=
\prod_p(1-p^{-s})^{-1},
\qquad \sigma>1.
\]

هذه الهوية هي الصورة التحليلية لوحيدية التحليل الأولي: الطرف الأيسر يجمع
على جميع الأعداد، والطرف الأيمن يبني كل عدد من قواه الأولية المحلية.

\section{المقلوب وسلسلة موبيوس}

\begin{corollary}[مقلوب زيتا في نصف مستوى التقارب المطلق]
\label{cor:mobius-reciprocal-zeta}
\resultid{ANT-COR-05-01}
\provedhere
إذا كان \(\Re(s)>1\)، فإن
\[
\frac1{\zeta(s)}
=
\sum_{n=1}^{\infty}\frac{\mu(n)}{n^s}
=
\prod_p(1-p^{-s}).
\]
وعلى وجه الخصوص، لا تنعدم \(\zeta(s)\) في النصف \(\Re(s)>1\).
\end{corollary}

\begin{proof}
لدينا في حلقة الالتفاف
\[
1*\mu=\varepsilon.
\]
وتتقارب سلسلتا \(1\) و\(\mu\) مطلقًا عندما \(\sigma>1\)، لأن
\(|\mu(n)|\leq1\). لذلك يعطي ضرب السلسلتين
\[
\zeta(s)
\sum_{n=1}^{\infty}\frac{\mu(n)}{n^s}
=
1.
\]
وهذا يثبت المقلوب وعدم الانعدام. أما المنتج، فينتج من
\[
\sum_{k=0}^{\infty}\frac{\mu(p^k)}{p^{ks}}
=1-p^{-s}.
\]
\end{proof}

\begin{remark}
عدم الانعدام المثبت هنا محصور في \(\sigma>1\). ولا يجوز نقله إلى حدود
أخرى قبل بناء الاستمرار التحليلي وأدوات المناطق الخالية من الأصفار.
\end{remark}

\section{المشتقة اللوغاريتمية وفون مانغولت}

\begin{theorem}[المشتقة اللوغاريتمية لدالة زيتا]
\label{thm:zeta-logarithmic-derivative-von-mangoldt}
\resultid{ANT-THM-05-05}
\provedhere
إذا كان \(\Re(s)>1\)، فإن
\[
-\frac{\zeta'(s)}{\zeta(s)}
=
\sum_{n=1}^{\infty}\frac{\Lambda(n)}{n^s}.
\]
\end{theorem}

\begin{proof}
يسمح التقارب المطلق بالاشتقاق حدًا حدًا:
\[
-\zeta'(s)
=
\sum_{n=1}^{\infty}\frac{\log n}{n^s}.
\]
ونضرب في
\[
\frac1{\zeta(s)}
=
\sum_{n=1}^{\infty}\frac{\mu(n)}{n^s}.
\]
وبحسب قاعدة ضرب سلاسل ديريشليه، تكون المعاملات هي
\[
(\mu*\log)(n).
\]
وقد أثبت الفصل الرابع أن
\[
\mu*\log=\Lambda.
\]
فتنتج الصيغة.
\end{proof}

\begin{remark}
تكشف المشتقة اللوغاريتمية العوامل الأولية: اللوغاريتم يحول الجداء إلى
مجموع، والاشتقاق يعطي أوزان \(\log p\) على القوى الأولية. لهذا تصبح
\(\Lambda\) المعامل الطبيعي عند دراسة توزيع الأوليات بواسطة زيتا.
\end{remark}

\section{هويات أساسية أخرى}

جميع الهويات الآتية صحيحة في أنصاف المستويات المذكورة، حيث تتقارب
السلاسل اللازمة مطلقًا.

\subsection{دالة القواسم}

من \(\tau=1*1\):
\[
\sum_{n=1}^{\infty}\frac{\tau(n)}{n^s}
=
\zeta(s)^2,
\qquad \sigma>1.
\]

\subsection{دالة أويلر}

من \(\varphi=\mu*\operatorname{id}\):
\[
\sum_{n=1}^{\infty}\frac{\varphi(n)}{n^s}
=
\frac{\zeta(s-1)}{\zeta(s)},
\qquad \sigma>2.
\]

فالتحويل الديريشلي لـ\(\operatorname{id}(n)=n\) هو
\[
\sum_{n=1}^{\infty}\frac{n}{n^s}=\zeta(s-1).
\]

\subsection{دالة ليوفيل}

من \(\lambda=\mu*1_{\square}\):
\[
\sum_{n=1}^{\infty}\frac{\lambda(n)}{n^s}
=
\frac{\zeta(2s)}{\zeta(s)},
\qquad \sigma>1.
\]

\subsection{دوال القواسم العامة}

من \(\sigma_\alpha=1*\operatorname{id}_\alpha\):
\[
\sum_{n=1}^{\infty}\frac{\sigma_\alpha(n)}{n^s}
=
\zeta(s)\zeta(s-\alpha),
\]
في النصف الآمن
\[
\sigma>\max(1,1+\Re(\alpha)).
\]

\section{مثال محلول: من هوية حسابية إلى هوية تحليلية}

نريد اشتقاق سلسلة \(\varphi\).

\begin{enumerate}
  \item الهوية الحسابية هي
  \[
  \varphi=\mu*\operatorname{id}.
  \]
  \item تحويل \(\mu\) هو
  \[
  D_\mu(s)=\frac1{\zeta(s)}
  \qquad(\sigma>1).
  \]
  \item تحويل \(\operatorname{id}\) هو
  \[
  D_{\operatorname{id}}(s)=\zeta(s-1)
  \qquad(\sigma>2).
  \]
  \item نختار تقاطع مجالي التقارب المطلق، أي \(\sigma>2\).
  \item نحول الالتفاف إلى ضرب:
  \[
  D_\varphi(s)
  =D_\mu(s)D_{\operatorname{id}}(s)
  =\frac{\zeta(s-1)}{\zeta(s)}.
  \]
\end{enumerate}

الخطوة الرابعة ليست تجميلًا تقنيًا؛ فهي التي تمنع استعمال الهوية خارج
المجال الذي يبرر ضرب المتسلسلتين وإعادة ترتيب الحدود.

\section{أسئلة تفسير البرهان}

\begin{enumerate}
  \item في مبرهنة نصف مستوى التقارب، أين استعملنا تقارب السلسلة عند
  \(s_0\)، ولماذا لا يكفي أن تكون حدودها فقط متجهة إلى الصفر؟
  \item لماذا أثبتنا التقارب المنتظم على المجموعات المدمجة بدل ادعائه على
  نصف المستوى كله؟
  \item في المنتج الأويلري، ما الدور المنفصل لكل من الضربية، ووحيدية
  التحليل الأولي، والتقارب المطلق؟
  \item لماذا تثبت هوية موبيوس عدم انعدام زيتا في \(\sigma>1\)، ولا تقول
  شيئًا مباشرًا عن الشريط الحرج؟
  \item في مبرهنة الفرادة، لماذا نختار أصغر معامل غير صفري؟
\end{enumerate}

\section{أخطاء شائعة}

\begin{enumerate}
  \item الخلط بين مجال تعريف السلسلة ومجال الاستمرار التحليلي للدالة.
  \item افتراض أن التقارب النقطي يكفي للاشتقاق حدًا حدًا.
  \item ضرب سلسلتين وإعادة ترتيب المجموع المزدوج دون تقارب مطلق أو بديل
  مبرهن.
  \item كتابة منتج أويلري لدالة غير ضربية.
  \item حذف شرط المجال من هويات صحيحة فقط في نصف مستوى معين.
  \item استنتاج عدم الانعدام من كون كل عامل محلي غير صفري من دون فحص
  تقارب الجداء اللانهائي.
  \item الخلط بين \(\sigma_c\) و\(\sigma_a\).
  \item استعمال \(-D'(s)/D(s)\) عند نقاط يمكن أن ينعدم فيها \(D(s)\).
\end{enumerate}

\section{تمارين}

\subsection{المستوى الأول: تثبيت المفاهيم}

\begin{enumerate}
  \item حدد نصفي مستوى التقارب والتقارب المطلق للسلسلة
  \[
  \sum_{n=1}^{\infty}\frac{(-1)^{n-1}}{n^s}.
  \]
  \item أثبت أن
  \[
  \sum_{n=1}^{\infty}\frac{1_{\square}(n)}{n^s}=\zeta(2s).
  \]
  \item اشتق حدًا حدًا سلسلة \(\zeta(s)\) في \(\sigma>1\).
  \item فسّر لماذا يؤدي \(A(x)=O(1)\) إلى التقارب في \(\sigma>0\).
\end{enumerate}

\subsection{المستوى الثاني: استعمال البنية}

\begin{enumerate}
  \item استعمل \(\tau_k=1*\cdots*1\) لإثبات
  \[
  \sum_{n=1}^{\infty}\frac{\tau_k(n)}{n^s}=\zeta(s)^k
  \qquad(\sigma>1).
  \]
  \item أثبت حاصل الضرب الأويلري لـ
  \[
  \sum_{n=1}^{\infty}\frac{\mu^2(n)}{n^s}
  \]
  ثم عبّر عنه بدلالة \(\zeta(s)\) و\(\zeta(2s)\).
  \item اشتق هوية دالة ليوفيل بطريقتين: من الالتفاف، ومن العوامل المحلية.
  \item أثبت صيغة \(\sigma_\alpha\) مع تحديد نصف مستوى مطلق التقارب.
\end{enumerate}

\subsection{المستوى الثالث: البرهان والتحليل}

\begin{enumerate}
  \item أكمل جميع تفاصيل التقدير المنتظم للذيل في برهان
  مبرهنة~\ref{thm:dirichlet-series-right-half-plane}.
  \item أثبت مبرهنة الفرادة باستخراج المعاملات واحدًا بعد آخر بواسطة نهايات
  على المحور الحقيقي.
  \item إذا كان \(A(x)=cx+O(x^\theta)\) مع \(\theta<1\)، فافصل من تمثيل
  أبيل الجزء الذي ينتج حدًا من شكل \(c\,s/(s-1)\)، وحدد مجال تكامل الباقي.
  \item ابحث عن مثال يحقق \(0<\sigma_a-\sigma_c<1\)، وبيّن بدقة أين
  تتقارب سلسلته وأين تتقارب مطلقًا.
\end{enumerate}

\section{مراجع الفصل ومسار التوسع}

للمعالجة الكلاسيكية لسلاسل ديريشليه والمنتجات الأويلرية يمكن الرجوع إلى
\cite{Apostol1976,Davenport2000,Titchmarsh1986}. أما الانتقال من هذه
المنطقة الآمنة إلى الاستمرار التحليلي والمعادلة الوظيفية ونظرية الأصفار،
فسيبدأ في الفصل التالي المخصص لدالة زيتا لريمان.

\section{خلاصة الفصل}

سلسلة ديريشليه تحول معاملات حسابية إلى دالة عقدية، ويحدد نمو المجاميع
الجزئية موضع تقاربها. ويحول التقارب المحلي المنتظم السلسلة إلى كائن
هولومورفي، بينما يحول التقارب المطلق الالتفاف إلى ضرب ويبرر المنتج
الأويلري. وعندما تكون المعاملات ضربية، تنقسم المعلومات إلى عوامل محلية عند
الأوليات؛ وعندما نأخذ المشتقة اللوغاريتمية لدالة زيتا تظهر دالة فون
مانغولت بوصفها الوزن الطبيعي للقوى الأولية.

هذه البنية هي نقطة الانطلاق الصحيحة لدراسة زيتا: لن نتعامل معها في الفصل
التالي كسلسلة منفردة، بل كدالة تجمع بين معاملاتها، ومنتجها الأويلري،
واستمرارها التحليلي، وأصفارها.
